\section{Discusión: qué podemos afirmar}\label{sec:discusion}

La Tabla~\ref{tab:resumen} reúne las pruebas. Ninguna causa del lugar pasa las tres, y la única explicación
compatible con todos los resultados es la forma de obtener la etiqueta en cada municipio. Esa explicación es la que
mejor encaja con el salto exacto en la frontera municipal y con el cambio de formato (dos decimales contra
redondeo).

\begin{table}[t]
\centering
\caption{Resumen de las hipótesis}
\label{tab:resumen}
\footnotesize
\begin{tabular}{@{}lccc@{}}
\toprule
Hipótesis & Difiere & Explica dentro & Se sostiene \\
 & entre A y B & del régimen & entre vecinos \\
\midrule
Clima & \si & \no & \no \\
Altitud & \no & \no & \no \\
Calendario y ciclo & \si & \no & \no \\
Suelo & \no & \no & \no \\
Ceniza volcánica & \si & \no & \no \\
Nivel oficial (SIAP) & \no & --- & \no \\
Origen de la etiqueta & \si & --- & \si \\
\bottomrule
\end{tabular}
\end{table}

\textbf{Lo que sí podemos afirmar.} La brecha de 1.44\,t/ha no se explica por el ambiente ni aparece en las cifras
oficiales. Está ligada a cómo se obtuvo el rendimiento en cada municipio, y la muestra del régimen A no representa el
nivel de su municipio.

\textbf{Lo que no podemos afirmar.} Con estos datos no se pueden separar dos causas:
\begin{itemize}
\item \emph{Sesgo de medición:} el rendimiento del régimen A se midió o se reportó de otra forma.
\item \emph{Sesgo de selección:} en los municipios A se muestrearon productores mejores que el promedio, por ejemplo
con contrato de compra, crédito o asistencia técnica, que sí obtienen alrededor de 4.5\,t/ha.
\end{itemize}
En el segundo caso sí habría una variable real, pero sería del productor y no del lugar. Tampoco podemos afirmar que
el régimen B esté sesgado. Su redondeo agrega ruido, de unas 0.14\,t/ha por parcela, pero su nivel queda cerca de lo
oficial.

\textbf{Implicaciones.} Para la calificación del reto no cambia nada: las 59 parcelas a predecir vienen de las mismas
fuentes, así que el modelo con régimen sigue siendo el adecuado. Para FIRA, el rendimiento en un lugar nuevo no se
puede estimar con el ambiente disponible, porque en estos datos el rendimiento depende más de quién siembra que de
dónde. Si el régimen A resulta ser un grupo de productores con acceso a crédito o a un esquema de contrato, ese
acceso se asociaría con casi el doble del rendimiento municipal oficial, y eso se conecta directamente con la
inclusión financiera.
